\documentclass[11pt]{article}
\usepackage[T1]{fontenc}
\usepackage{lmodern,amsmath,amssymb,amsthm,mathtools}
\usepackage[margin=1.05in]{geometry}
\usepackage[hidelinks]{hyperref}
\newtheorem{theorem}{Theorem}[section]
\newtheorem{lemma}[theorem]{Lemma}
\newtheorem{proposition}[theorem]{Proposition}
\newtheorem{corollary}[theorem]{Corollary}
\newcommand{\R}{\mathbb R}\newcommand{\N}{\mathbb N}
\newcommand{\E}{\mathbb E}\newcommand{\Ent}{\operatorname{Ent}}
\newcommand{\Dom}{\operatorname{Dom}}
\newcommand{\code}[1]{{\small\nolinkurl{#1}}}
\title{Poisson--Gaussian transport:\\closed graphs, coherent recovery and nonlinear limits}
\author{Project Stardust}
\date{Reviewed formalization companion, October 29, 2026}
\begin{document}
\maketitle
\begin{abstract}
The Charlier--Hermite correspondence identifies Poisson and Gaussian
Hilbert spaces and their complete annihilation and creation graphs.
We reconstruct the basis, maximal domains and graph-core arguments,
then transport actual Poisson changes of intensity into a strongly
continuous Gaussian displacement group with an explicit weighted-shift
formula on the whole $L^2$ space. The correspondence preserves distances
between square-root likelihoods and their derivative energies, but it
does not preserve positivity, ordinary multiplication, or entropy.
Coherent vectors give quantitative entropy distortion and exclude a
bounded realization of any nonzero displacement on $L^4$.
A compact creation inverse and a sharp shifted Poincar\'e estimate hold on the original maximal domains. A compensated-count example separates the small- and large-intensity $L^{4/3}$ scales.
Positive pointwise jump curvature coexists with this entropy failure,
and the missing diffusion chain rule is exhibited at one atom.
The failure of a quadratic Poisson logarithmic Sobolev inequality already
occurs on finite-support functions; we prove this by controlling the
cutoff boundary in the original difference graph. These are classical
mechanisms, consolidated as a companion to frozen verified constructions.
No new Fock-space correspondence or logarithmic Sobolev failure is claimed.
\end{abstract}

\section{What is transported}
Fix $\lambda>0$. Write
\[
 p_\lambda(k)=e^{-\lambda}\frac{\lambda^k}{k!},\quad k\in\N,
 \qquad d\gamma(x)=\frac{e^{-x^2/2}}{\sqrt{2\pi}}\,dx,
\]
and use real Hilbert spaces $H_\lambda=L^2(p_\lambda)$ and
$H_G=L^2(\gamma)$. Since every Poisson atom has positive mass, an element
of $H_\lambda$ has a uniquely determined value at every integer.
Gaussian functions are equivalence classes modulo $\gamma$-null sets.
The scalar normalization of the Poisson annihilation operator is
\begin{equation}\label{eq:A}
 A_\lambda f(k)=\sqrt\lambda\,[f(k+1)-f(k)].
\end{equation}
Its domain consists of precisely those $f\in H_\lambda$ for which
the displayed sequence belongs to $H_\lambda$. We will construct a
unitary $U_\lambda:H_\lambda\to H_G$ and prove, on full domains,
\[
 U_\lambda A_\lambda=D U_\lambda,\qquad
 U_\lambda A_\lambda^*=\delta U_\lambda,
 \quad D=\partial_x,\quad\delta=x-\partial_x.
\]
This is an identification of Hilbert spaces with closed operators.
It supplies no identification of their ordinary products or laws.

The correspondence is classical. Vershik and Tsilevich
\cite[Sections 1.3 and 3.1]{VT} give the Poisson--Gaussian isometry,
including its Hermite--Charlier specialization, and explicitly discuss
its failure to be positive or multiplicative. The quadratic Poisson
entropy obstruction is also classical; the introduction and Theorem 1.1
of Aryan--L\'opez-Rivera--Shenfeld \cite{ALRS} place it beside Wu's
modified inequality. Our contribution is a complete exposition tied to
the original verified graphs, gauges and counterexamples, with these
antecedents distinguished from claims of priority.

\section{Complete bases and the unitary correspondence}
Define the Charlier and probabilists' Hermite polynomials by
\begin{align}
 e^{-\lambda t}(1+t)^k&=\sum_{n\geq0} C_n^\lambda(k)\frac{t^n}{n!},
 & e^{tx-t^2/2}&=\sum_{n\geq0}\operatorname{He}_n(x)\frac{t^n}{n!}.
 \label{eq:generating}
\end{align}
Thus $C_0=1$, $C_1=k-\lambda$, and
$C_2=(k-\lambda)^2-(k-\lambda)-\lambda$. Put
\[
 q_n^\lambda=\frac{C_n^\lambda}{\sqrt{n!\lambda^n}},\qquad
 h_n=\frac{\operatorname{He}_n}{\sqrt{n!}}.
\]

\begin{lemma}\label{lem:basis}
The sequences $(q_n^\lambda)$ and $(h_n)$ are complete orthonormal
bases of $H_\lambda$ and $H_G$, respectively. There is exactly one
unitary $U_\lambda$ sending every $q_n^\lambda$ to $h_n$.
\end{lemma}
\begin{proof}
For $N$ of law $p_\lambda$, the expectation of the product of the first
generating functions at $s,t$ is
\[
 e^{-\lambda(s+t)}\E[(1+s)^N(1+t)^N]=e^{\lambda st}.
\]
The analogous Gaussian expectation is $e^{st}$. Differentiating at zero
gives orthogonality with squared norms $n!\lambda^n$ and $n!$.
All differentiations are justified locally by exponential integrability.
Both polynomial sequences are monic and have degree $n$, so their
algebraic spans are the polynomial spaces.

To prove Poisson completeness, suppose $f\in H_\lambda$ is orthogonal
to all polynomials. The series
\[
 F(z)=\sum_{k\geq0} f(k)p_\lambda(k)z^k
\]
is entire: on $|z|\leq R$ its absolute sum is bounded by
$\|f\|_2\exp\{\lambda(R^2-1)/2\}$, and its tails tend to zero
uniformly by Cauchy--Schwarz. Every derivative at $z=1$ vanishes,
because falling factorials are polynomials. Thus $F=0$, and every
coefficient $f(k)p_\lambda(k)$ is zero.

For Gaussian completeness, if $f$ is orthogonal to every polynomial,
the entire function $L(z)=\int f(x)e^{zx}\,d\gamma(x)$ has all
derivatives at zero equal to zero. Entirety and differentiation follow
by Cauchy--Schwarz and Gaussian exponential moments, uniformly on
compact sets. Hence $L=0$. The finite signed measure $f\,d\gamma$
has zero Fourier transform and is zero, so $f=0$. For completeness,
Fourier uniqueness here follows by convolution with a Gaussian
approximate identity: Fourier inversion gives a zero convolution at
every positive variance, and the approximate identity converges weakly
to the original finite measure. Finally extend the isometry between
the two orthonormal bases by Hilbert completion; density gives both
surjectivity and uniqueness.
\end{proof}

For $f_n=\langle q_n^\lambda,f\rangle$, the construction gives
\begin{equation}\label{eq:U}
 U_\lambda f=\sum_{n\geq0}f_n h_n,\qquad
 \|f\|_2^2=\sum_{n\geq0}|f_n|^2.
\end{equation}
Every series in this paper that represents an arbitrary $L^2$ function
is understood in its Hilbert norm, rather than pointwise.

\begin{proposition}\label{prop:product}
The unitary $U_\lambda$ is neither positivity preserving nor
multiplicative for ordinary products, even on polynomials.
\end{proposition}
\begin{proof}
It sends the nonnegative function $k$ to $\lambda+\sqrt\lambda x$,
which is negative on a set of positive Gaussian measure. For
$r(k)=(k-\lambda)/\sqrt\lambda$, we have $U_\lambda r=x$ whereas
the identity for $C_2$ gives
\[
 U_\lambda(r^2)=x^2+\frac{x}{\sqrt\lambda}\ne (U_\lambda r)^2.
\]
All functions and products used here belong to $L^2$.
\end{proof}

\section{The maximal graphs, including their endpoints}
Define the maximal creation expression by
\begin{equation}\label{eq:creation}
 A_\lambda^*g(k)=\frac{k}{\sqrt\lambda}g(k-1)-\sqrt\lambda g(k),
\end{equation}
where the first term is zero at $k=0$. Its domain is the set of
$g\in H_\lambda$ for which this entire expression belongs to
$H_\lambda$; separate square integrability of its two terms is not
required. Both maximal graphs are closed. Indeed, convergence in
$H_\lambda$ implies convergence at every fixed atom, since
$|f(k)|\leq p_\lambda(k)^{-1/2}\|f\|_2$. One may therefore pass to
the pointwise expressions in any convergent graph sequence.

\begin{theorem}\label{thm:graphs}
For $f,g\in H_\lambda$ with coefficients $(f_n),(g_n)$,
\begin{align}
 A_\lambda f=g
 &\ \Longleftrightarrow\quad
 g_n=\sqrt{n+1}\,f_{n+1}\quad(n\geq0),\label{eq:lower}\\
 A_\lambda^*g=f
 &\ \Longleftrightarrow\quad
 f_0=0,\quad f_{n+1}=\sqrt{n+1}\,g_n\quad(n\geq0).
 \label{eq:raise}
\end{align}
Consequently
\begin{align}
 \Dom A_\lambda&=\left\{f:\sum_{n\geq0}n|f_n|^2<\infty\right\},
 &\|A_\lambda f\|_2^2&=\sum_{n\geq0}n|f_n|^2,\label{eq:Adomain}\\
 \Dom A_\lambda^*&=\left\{g:\sum_{n\geq0}(n+1)|g_n|^2<\infty\right\},
 &\|A_\lambda^*g\|_2^2&=\sum_{n\geq0}(n+1)|g_n|^2.
 \label{eq:Astar}
\end{align}
These operators are Hilbert adjoints of one another, and finite
Charlier expansions are graph cores for both.
\end{theorem}
\begin{proof}
The generating function gives
$\Delta C_{n+1}^\lambda=(n+1)C_n^\lambda$ and hence
$A_\lambda q_{n+1}^\lambda=\sqrt{n+1}q_n^\lambda$.
Summation by parts uses exactly
$(k+1)p_\lambda(k+1)=\lambda p_\lambda(k)$, including the
zero endpoint in \eqref{eq:creation}. Against a polynomial test it yields
\[
 \langle A_\lambda f,v\rangle=\langle f,A_\lambda^*v\rangle,
 \qquad
 \langle A_\lambda^*g,v\rangle=\langle g,A_\lambda v\rangle
\]
for the respective maximal-domain inputs. The reindexed sums are
absolutely convergent: each reduces to a product of an $L^2$ input
and a polynomial in $L^2(p_\lambda)$. Thus no boundary term at infinity
is discarded. Applying the first identity to $q_n^\lambda$ and the
second to $q_{n+1}^\lambda$ gives the coefficient rules. The test $1$
also gives $f_0=0$ in the creation rule.

Conversely, assume the coefficient rule \eqref{eq:lower}. Truncate
$f$ at indices $0,\ldots,M$ and $g$ at $0,\ldots,M-1$.
The finite ladder identities give the annihilation graph at every
truncation; both components converge by completeness. Closedness proves
the full graph relation. For \eqref{eq:raise}, truncate $g$ at
$0,\ldots,M-1$ and $f$ at $0,\ldots,M$, and use the same argument.
This also proves the stated domain formulas and graph-core property.

In coefficient coordinates the annihilation operator is the weighted
backward shift $(f_n)\mapsto(\sqrt{n+1}f_{n+1})$.
Its Hilbert adjoint is precisely the weighted forward shift in
\eqref{eq:raise}: finite sequences determine the adjoint coordinates,
and their square summability is necessary and sufficient for a bounded
functional on $\ell^2$. This proves actual adjointness on the full
domains, rather than just the polynomial integration-by-parts identity.
\end{proof}

\begin{lemma}\label{lem:finitecore}
Finite-support sequences are also a common graph core for
$A_\lambda$ and $A_\lambda^*$.
\end{lemma}
\begin{proof}
Write $f^{[M]}(k)=f(k)\mathbf1_{\{k<M\}}$. For
$f\in\Dom A_\lambda$, the shifted sequence
$f(k+1)=f(k)+A_\lambda f(k)/\sqrt\lambda$ is square integrable.
Reindexing shows $\sum_k k p_\lambda(k)f(k)^2<\infty$. Thus
\[
 \|A_\lambda f^{[M]}-A_\lambda f\|_2^2
 =M p_\lambda(M)f(M)^2+
 \sum_{k\geq M}p_\lambda(k)|A_\lambda f(k)|^2\longrightarrow0.
\]
For $g\in\Dom A_\lambda^*$, direct evaluation at the single boundary
atom gives instead
\[
 \|A_\lambda^*g^{[M]}-A_\lambda^*g\|_2^2
 =\lambda p_\lambda(M)g(M)^2+
 \sum_{k>M}p_\lambda(k)|A_\lambda^*g(k)|^2\longrightarrow0.
\]
Both scalar cutoffs converge in $H_\lambda$, and their operator
expressions have finite support. These identities prove graph density
in both domains, including the boundary contributions.
\end{proof}

\begin{theorem}\label{thm:transport}
Let $D$ be the maximal weak derivative in $L^2(\gamma)$ and let
$\delta g=xg-g'$ in the distributional sense, with maximal $L^2$
domain. Then
\[
 (f,g)\in\operatorname{graph} A_\lambda
 \iff (U_\lambda f,U_\lambda g)\in\operatorname{graph} D,
\]
and the same equivalence holds with $A_\lambda^*$ and $\delta$.
\end{theorem}
\begin{proof}
The Gaussian generating function gives $Dh_{n+1}=\sqrt{n+1}h_n$
and $\delta h_n=\sqrt{n+1}h_{n+1}$. The weak graphs are closed:
Gaussian $L^2$ convergence implies local Lebesgue $L^1$ convergence,
so the distributional equations pass to limits. For a maximal weak
graph, Gaussian integration by parts against $h_n\chi_R$, where
$\chi_R$ is a smooth cutoff equal to one on $[-R,R]$, gives the same
coefficient rules as \eqref{eq:lower}--\eqref{eq:raise} as $R\to\infty$.
The polynomial terms converge in $L^2(\gamma)$; the cutoff-derivative
term tends to zero by Cauchy--Schwarz and Gaussian polynomial tails.
This argument applies to $xg-g'$ as one distribution, without requiring
$xg$ and $g'$ separately to belong to $L^2$. Conversely, finite Hermite
truncations and closedness give both maximal weak graphs from their
coefficient rules. Their domains therefore coincide exactly with the
two weighted sequence domains in Theorem~\ref{thm:graphs}. Equation
\eqref{eq:U} now proves both equivalences.
\end{proof}

In the frozen formalization the original difference is $\Delta$, with
field space $L^2(\lambda p_\lambda)$. Multiplication of a field by
$\sqrt\lambda$ identifies that space with $H_\lambda$ and produces
\eqref{eq:A}. This normalization is essential for the graph transport.

\section{Likelihood roots and the physical change of intensity}
For $\theta\in\R$ put $\lambda_\theta=\lambda e^\theta$ and
\begin{equation}\label{eq:roots}
 r_{\lambda,\theta}(k)=
 \exp\!\left\{\frac{\theta k}{2}-\frac{\lambda(e^\theta-1)}2\right\},
 \quad a=\sqrt\lambda(e^{\theta/2}-1),\quad
 g_a(x)=e^{ax-a^2}.
\end{equation}
The square of $r_{\lambda,\theta}$ is the actual likelihood ratio
$p_{\lambda_\theta}/p_\lambda$; the square of $g_a$ is the
likelihood ratio of $N(2a,1)$ to $N(0,1)$. Both roots have norm one.

\begin{proposition}\label{prop:root}
One has $U_\lambda r_{\lambda,\theta}=g_a$. For any two tilt parameters,
\begin{align}
 \langle r_{\lambda,\theta},r_{\lambda,\eta}\rangle
 &=\exp\!\left\{-\tfrac\lambda2(e^{\theta/2}-e^{\eta/2})^2\right\},
 \label{eq:hellinger}\\
 \|r_{\lambda,\theta}-r_{\lambda,\eta}\|_2^2
 &=2\left[1-\exp\!\left\{-\tfrac\lambda2
 (e^{\theta/2}-e^{\eta/2})^2\right\}\right].\nonumber
\end{align}
Thus the correspondence preserves the exact Hellinger distances within
this family, under the square-root-intensity coordinate.
\end{proposition}
\begin{proof}
Write $s=e^{\theta/2}$, so $a=\sqrt\lambda(s-1)$.
Poisson summation of the first generating function gives
\[
 \E\big[r_{\lambda,\theta}(N)e^{-\lambda t}(1+t)^N\big]
 =e^{-a^2/2}e^{\sqrt\lambda a t}.
\]
Therefore the $n$th normalized Charlier coefficient of the root is
$e^{-a^2/2}a^n/\sqrt{n!}$. These are precisely the normalized Hermite
coefficients of $g_a$, by the second generating function. Coefficient
equality in the complete bases proves the identity. Directly summing
the product of the two Poisson roots gives \eqref{eq:hellinger}.
The norm formula follows because both roots have norm one. Alternatively,
the Gaussian identity $\langle g_a,g_b\rangle=e^{-(a-b)^2/2}$ gives
the same calculation after transport.
\end{proof}

The half-density change of measure is the actual unitary
\begin{equation}\label{eq:gauge}
 V_{\lambda,\theta}:H_{\lambda_\theta}\longrightarrow H_\lambda,
 \qquad V_{\lambda,\theta}f=r_{\lambda,\theta}f.
\end{equation}
Indeed, $\|Vf\|_{H_\lambda}^2=\|f\|_{H_{\lambda_\theta}}^2$,
and division by the strictly positive root is the inverse.

\begin{proposition}\label{prop:gauge}
On the complete maximal graphs,
\begin{align}
 A_\lambda Vf&=V A_{\lambda_\theta}f+aVf,
 &\Dom A_\lambda&=V\Dom A_{\lambda_\theta},\label{eq:gaugeA}\\
 A_\lambda^*Vf&=V A_{\lambda_\theta}^*f+aVf,
 &\Dom A_\lambda^*&=V\Dom A_{\lambda_\theta}^*.
 \label{eq:gaugecreation}
\end{align}
\end{proposition}
\begin{proof}
Use $r(k+1)=e^{\theta/2}r(k)$ in \eqref{eq:A} and
$r(k-1)=e^{-\theta/2}r(k)$ in \eqref{eq:creation}.
The two pointwise identities follow by expansion, including $k=0$.
Since $V$ is unitary, the extra term $aVf$ is always in $H_\lambda$.
Consequently either side has an $L^2$ operator output exactly when the
other does. This proves both domain equalities and both graph identities.
\end{proof}

\section{A group on the entire Gaussian space}
Define
\begin{equation}\label{eq:W}
 W_a=U_\lambda V_{\lambda,\theta}U_{\lambda_\theta}^{-1},
 \qquad a=\sqrt\lambda(e^{\theta/2}-1).
\end{equation}
Every $a\in\R$ is attainable: choose $b>|a|$, take $\lambda=b^2$
and $\theta=2\log((b+a)/b)$. The following result proves that
\eqref{eq:W} is independent of this choice.

\begin{theorem}\label{thm:physical}
For every $f\in H_G$, the whole-space formula is
\begin{equation}\label{eq:physical}
 (W_af)(x)=e^{ax-a^2}f(x-2a)\qquad\text{for almost every }x.
\end{equation}
These operators form a strongly continuous unitary group:
$W_aW_b=W_{a+b}$, $W_0=I$. They satisfy, on exact domains,
\begin{equation}\label{eq:commute}
 DW_af=W_aDf+aW_af,\qquad
 \delta W_af=W_a\delta f+aW_af.
\end{equation}
\end{theorem}
\begin{proof}
First the right side of \eqref{eq:physical} defines a unitary $T_a$.
The change of variables $y=x-2a$ gives
\[
 e^{2ax-2a^2}e^{-x^2/2}=e^{-(x-2a)^2/2},\qquad
 \|T_af\|_2^2=\|f\|_2^2.
\]
All Gaussian null sets are Lebesgue null sets, and translations preserve
them; hence the formula is well defined on equivalence classes.
The same formula gives $T_a^{-1}=T_{-a}$.

We identify $T_a$ with the unitary in \eqref{eq:W} without assuming a
polynomial rule extends to a nonlinear expression. For $s=\sqrt{
\lambda_\theta}$ and any $b>-s$, choose $\eta$ so that
$b=s(e^{\eta/2}-1)$. Proposition~\ref{prop:root} and the exact cocycle
\[
 r_{\lambda,\theta}\,r_{\lambda_\theta,\eta}
 =r_{\lambda,\theta+\eta}
\]
give $W_ag_b=g_{a+b}$. Direct substitution gives the same identity
for $T_a$. The vectors $g_b$ with $b$ in any nonempty open interval
have dense linear span in $H_G$: orthogonality to them makes the
entire function $\int f(x)e^{zx}\,d\gamma(x)$ vanish on that interval;
the identity theorem and Fourier uniqueness then give $f=0$.
Thus the bounded operators $W_a,T_a$ agree on a dense set and everywhere.

Multiplying the two explicit weights in $T_aT_b$ gives
$e^{(a+b)x-(a+b)^2}$ and translation by $2(a+b)$, proving the group
law and $W_0=I$. Also
\[
 \|W_ag_b-g_b\|_2^2=2(1-e^{-a^2/2})\longrightarrow0.
\]
Finite coherent linear combinations are dense, and the $W_a$ are
isometries, so this proves strong continuity at zero for every $f$.
The group law proves it at every parameter.
Finally \eqref{eq:commute} follows from
Theorem~\ref{thm:transport} and Proposition~\ref{prop:gauge}, including
domain equivalence. It can also be checked by distributional
differentiation of \eqref{eq:physical}; the inverse shift then gives
the converse domain inclusion.
\end{proof}

Strong continuity does not mean continuity in operator norm.
For $a>0$, let $f_a$ be the indicator of $(0,a)$ divided by its
Gaussian $L^2$ norm. Its displaced image is supported on $(2a,3a)$.
The supports are disjoint, while both vectors have norm one. Consequently
\begin{equation}\label{eq:normdiscontinuity}
 \|W_a-I\|_{H_G\to H_G}\geq\|(W_a-I)f_a\|_2=\sqrt2
 \qquad(a>0).
\end{equation}
This lower bound suffices to exclude operator-norm continuity at zero;
no optimal operator-norm value is needed.

\section{Quartic and entropy obstructions}
\begin{theorem}\label{thm:L4}
For every fixed $a\ne0$, no bounded operator on $L^4(\gamma)$ can
agree with $W_a$ on all coherent vectors $g_b$, $b\in\R$.
In particular, the $L^2$ unitary does not admit a bounded $L^4$
realization with its coherent action.
\end{theorem}
\begin{proof}
Gaussian exponential moments give $\|g_b\|_4=e^{b^2}$, while
Theorem~\ref{thm:physical} gives $W_ag_b=g_{a+b}$. Thus
\[
 \frac{\|W_ag_b\|_4}{\|g_b\|_4}
 =e^{2ab+a^2}.
\]
This diverges as $b\to+\infty$ if $a>0$, or $b\to-\infty$ if $a<0$.
It contradicts any fixed bounded-operator norm. The argument uses only
coherent functions, which belong to both $L^2$ and $L^4$.
\end{proof}

For a probability measure $\mu$ write
\[
 \Ent_\mu(F)=\int F\log F\,d\mu
 -\left(\int F\,d\mu\right)\log\left(\int F\,d\mu\right),
 \qquad 0\log0=0,
\]
whenever the terms are finite. The roots in \eqref{eq:roots} have
normalized squares, so their square entropies are relative entropies.

\begin{theorem}\label{thm:entropy}
For $a=\sqrt\lambda(e^{\theta/2}-1)$ one has
\begin{align}
 \Ent_{p_\lambda}(r_{\lambda,\theta}^2)
 &=\lambda(\theta e^\theta-e^\theta+1),\label{eq:Poissonentropy}\\
 \Ent_\gamma(g_a^2)&=2a^2,
 &\|A_\lambda r_{\lambda,\theta}\|_2^2
 &=\|Dg_a\|_2^2=a^2.\label{eq:Gaussianentropy}
\end{align}
For every $K\geq0$ there is a positive $\theta$ such that
\begin{equation}\label{eq:distortion}
 K\Ent_\gamma(g_a^2)<\Ent_{p_\lambda}(r_{\lambda,\theta}^2).
\end{equation}
Hence exact root distances and derivative energies do not give a
uniform upper comparison of these ordinary entropies.
\end{theorem}
\begin{proof}
The log of the Poisson likelihood is
$\theta k-\lambda(e^\theta-1)$. Its expectation under the tilted
Poisson law uses $\E_{p_{\lambda e^\theta}}k=\lambda e^\theta$,
giving \eqref{eq:Poissonentropy}. All absolute logarithmic moments
are finite by Poisson exponential integrability. For the Gaussian root,
$g_a^2\gamma=N(2a,1)$ and $\log g_a^2=2ax-2a^2$, giving entropy
$2a^2$. The identities $A_\lambda r=ar$ and $Dg_a=ag_a$
give the energies because both roots have norm one.

Choose $\theta=2K+2$. Since $\theta>0$,
\[
 \Ent_\gamma(g_a^2)=2\lambda(e^{\theta/2}-1)^2
 \leq2\lambda e^\theta,
\]
whereas
$\Ent_{p_\lambda}(r^2)=\lambda[(\theta-1)e^\theta+1]
>2K\lambda e^\theta$. This proves \eqref{eq:distortion}.
\end{proof}

\begin{corollary}\label{cor:core}
For every $C\geq0$ there is a real finite-support function $f$ on
$\N$ such that
\begin{equation}\label{eq:corefailure}
 C\lambda\sum_{k\geq0}p_\lambda(k)[f(k+1)-f(k)]^2
 <\Ent_{p_\lambda}(f^2).
\end{equation}
Such an $f$ belongs to the original difference domain and to the
original creation domain. Thus failure of a quadratic Poisson
logarithmic Sobolev estimate already occurs on this common core.
\end{corollary}
\begin{proof}
By Theorem~\ref{thm:entropy}, choose $\theta>0$ with
$\Ent_{p_\lambda}(r^2)>C a^2$, where $r=r_{\lambda,\theta}$.
Set $f_M(k)=r(k)\mathbf1_{\{k<M\}}$ for $M\geq1$.
The $L^2$ tail tends to zero. The difference between
$\Delta f_M$ and $\Delta r$ is zero below $M-1$, equals $-r(M)$
at $M-1$, and equals $-\Delta r(k)$ for $k\geq M$.
Consequently
\begin{equation}\label{eq:boundary}
 \|A_\lambda f_M-A_\lambda r\|_2^2
 =\lambda p_\lambda(M-1)r(M)^2
  +\lambda\sum_{k\geq M}p_\lambda(k)[\Delta r(k)]^2
 \longrightarrow0.
\end{equation}
The second term is a tail of a convergent series. The first is a
constant times $\lambda^M e^{\theta M}/(M-1)!$, which tends to zero
by the ratio test. This checks the otherwise significant cutoff
boundary in the actual maximal graph.

Since $r(k)^2|\log r(k)^2|$ is bounded by a constant times
$e^{\theta k}(1+k)$, it is Poisson integrable. Therefore
\[
 \sum p_\lambda f_M^2\log f_M^2\longrightarrow
 \sum p_\lambda r^2\log r^2,
 \qquad \sum p_\lambda f_M^2\longrightarrow1.
\]
Continuity of $x\log x$ at $1$ proves
$\Ent(f_M^2)\to\Ent(r^2)$. Equation \eqref{eq:boundary} gives
$\|A_\lambda f_M\|_2^2\to a^2$. The strict entropy-energy margin
thus persists for all sufficiently large $M$, giving
\eqref{eq:corefailure}. Finally both difference and creation expressions
of a finite-support sequence have finite support, so both original
domain assertions hold directly.
\end{proof}

This obstruction concerns quadratic difference energy. It does not
contradict the modified Poisson inequality with jump Bregman terms.
For clarity, let $F=r_{\lambda,\theta}^2$. Its jump is
$F(k+1)=e^\theta F(k)$, and hence
\[
 \lambda\sum_k p_\lambda(k)
 \left[F(k+1)\log\frac{F(k+1)}{F(k)}-F(k+1)+F(k)\right]
 =\lambda(\theta e^\theta-e^\theta+1)=\Ent_{p_\lambda}(F).
\]
The same exponentials therefore attain equality in the classical
modified estimate \cite{ALRS}, while excluding every fixed quadratic
constant. Ordinary product structure determines which energy is relevant.

\section{Positive quadratic curvature and the missing chain rule}
The same original graphs identify the immigration--death expression
\begin{equation}\label{eq:jumpL}
 Lf(k)=\lambda[f(k+1)-f(k)]+k[f(k-1)-f(k)]
       =-A_\lambda^*A_\lambda f(k).
\end{equation}
At zero the death term is zero. The last equality is pointwise whenever
the two original graph memberships hold, and in particular on finite
support. For arbitrary functions on $\N$ define, pointwise,
\[
 \Gamma(f,g)=\tfrac12[L(fg)-fLg-gLf],\quad
 \Gamma(f)=\Gamma(f,f),\quad
 \Gamma_2(f)=\tfrac12 L\Gamma(f)-\Gamma(f,Lf).
\]
Every expression uses only finitely many neighboring values. There is
no summability assumption in a pointwise curvature statement.
This quadratic jump calculus is different from diffusion chain calculus;
Weber--Zacher \cite[Example 5.13]{WZ} already use the Poisson process to
separate positive Bakry--\'Emery curvature from their entropy-adapted
discrete curvature. We make no new curvature claim.

\begin{theorem}\label{thm:jumpcurvature}
For every real function $f$ on $\N$,
\begin{equation}\label{eq:jumpGamma}
 \Gamma(f)(k)=\frac\lambda2[f(k+1)-f(k)]^2
       +\frac k2[f(k)-f(k-1)]^2,
 \qquad \Gamma_2(f)(k)\geq\tfrac12\Gamma(f)(k).
\end{equation}
The constant $1/2$ is the largest constant valid uniformly in $k$ and
$f$, even if the tested functions are required to have finite support.
Nevertheless, for every $C\geq0$ there is a finite-support $f$ with
\begin{equation}\label{eq:curvatureentropyfailure}
 C\int\Gamma(f)\,dp_\lambda<\Ent_{p_\lambda}(f^2).
\end{equation}
\end{theorem}
\begin{proof}
Expanding $L(fg)-fLg-gLf$ gives
\[
 \Gamma(f,g)(k)=\tfrac\lambda2\Delta f(k)\Delta g(k)
 +\tfrac k2[f(k)-f(k-1)][g(k)-g(k-1)],
\]
and hence the first formula. To see all terms in the second assertion,
put $u=f(k+1)-f(k)$, $v=f(k)-f(k-1)$ and
\[
 A=f(k+2)-2f(k+1)+f(k),\qquad
 B=f(k)-2f(k-1)+f(k-2).
\]
For $k\geq2$, the neighboring first differences are $u+A$ and $v-B$,
so
\begin{align*}
 \Gamma(f)(k+1)&=\tfrac\lambda2(u+A)^2+\tfrac{k+1}2u^2,\\
 \Gamma(f)(k-1)&=\tfrac\lambda2v^2+\tfrac{k-1}2(v-B)^2,\\
 \Delta Lf(k)&=\lambda A-u-k(u-v),\\
 Lf(k)-Lf(k-1)&=\lambda(u-v)-v-(k-1)B.
\end{align*}
Substitute these four identities into the definition of $\Gamma_2$.
The linear cross terms in $A,B$ cancel; the remaining $u,v$ cross
term is absorbed into $(u-v)^2$. The result is the exact decomposition
\begin{equation}\label{eq:jumpGamma2}
 \begin{split}
 \Gamma_2(f)(k)={}&\frac{\lambda^2}4 A^2+
       \frac{k(k-1)}4 B^2+\frac{\lambda k}2(u-v)^2\\
       &+\frac{3\lambda}4u^2+\frac k4v^2.
 \end{split}
\end{equation}
At $k=0$ the terms with $k$ vanish, and at $k=1$ the $B$ term
vanishes; direct evaluation of the defining expressions gives the same
formula at both endpoints. Negative-index values can be set to zero,
since their terms have zero coefficient. Subtracting
$\tfrac12\Gamma(f)$ leaves the three nonnegative square terms and
$\lambda u^2/2$, proving the bound.

For the count function $f(k)=k$, one obtains
\[
 \Gamma(f)(k)=\frac{\lambda+k}{2},\qquad
 \Gamma_2(f)(k)=\frac{3\lambda+k}{4},\qquad
 \frac{\Gamma_2(f)(k)}{\Gamma(f)(k)}\longrightarrow\frac12.
\]
Any proposed larger constant fails at some finite $k$. Replace this
count function by a finite-support function agreeing with it on
$\{0,\ldots,k+2\}$; both pointwise expressions at that $k$ are unchanged.
This proves the finite-support sharpness assertion as an elementary
local deduction from the verified count witness.

Finally detailed balance $k p_\lambda(k)=\lambda p_\lambda(k-1)$
gives, for finite-support $f$,
\[
 \int\Gamma(f)\,dp_\lambda
 =\lambda\sum_{k\geq0}p_\lambda(k)[\Delta f(k)]^2
 =\|A_\lambda f\|_2^2.
\]
Indeed the death-square sum reindexes to the birth-square sum, each
contributing half of the last energy. Corollary~\ref{cor:core} now gives
\eqref{eq:curvatureentropyfailure}, for the same finite-support domain
on which the curvature estimate holds everywhere.
\end{proof}

The missing diffusion step can be seen at one atom. For the count
function $\iota(k)=k$ and $\Phi(x)=e^x$,
\[
 \Gamma(\Phi\circ\iota)(0)=\frac\lambda2(e-1)^2
 >\frac\lambda2=\Phi'(\iota(0))^2\Gamma(\iota)(0).
\]
The strict inequality uses $e>2$. Since $p_\lambda(0)>0$, even an
almost-everywhere version of that diffusion chain rule is false.
The source energy $E_f(k)=\lambda[\Delta f(k)]^2$ is also different
from $\Gamma(f)(k)$ pointwise: for $f=\iota$ at zero the values are
$\lambda$ and $\lambda/2$. Their integrals agree on the core, which
does not make their pointwise identities interchangeable.

There is even a unit lower Bochner bound for this different forward
energy. Direct evaluation gives
$\Delta Lf=L\Delta f-\Delta f$. Applying
$\tfrac12 L(h^2)=hLh+\Gamma(h)$ with $h=\Delta f$ therefore yields
\[
 \tfrac12 L E_f-\lambda\Delta f\,\Delta Lf
 =E_f+\lambda\Gamma(\Delta f)\geq E_f.
\]
This identity and the sharp $1/2$ bound for the generator's
$\Gamma$ are both valid. Neither supplies the missing chain rule or
the false quadratic entropy estimate. The original graph equivalence
with Gaussian operators must be used together with the product
obstruction in Proposition~\ref{prop:product}.



\section{Full-domain recovery and a displaced vacuum}
The preceding equivalence also identifies the exact range of creation.
We give the scalar normalization explicitly, so that no weighted-field
factor is hidden in the estimates. Write $p_n=\varphi_n^\lambda$ for the
orthonormal Charlier basis of Section 2, with $p_0=1$, and set
\[
 C_\lambda f=\sum_{n\geq0}\frac{\langle p_{n+1},f\rangle}
 {\sqrt{n+1}}p_n,\qquad f^\circ=f-\E_\lambda f.
\]
\begin{proposition}\label{prop:recovery}
The bounded compact operator $C_\lambda:H_\lambda\to H_\lambda$
satisfies, for every $f\in H_\lambda$,
\[
 C_\lambda f\in\Dom(A_\lambda^*),\quad
 A_\lambda^*C_\lambda f=f^\circ,\quad
 \|C_\lambda f\|_2^2=
 \sum_{q\geq1}\frac{|\langle p_q,f\rangle|^2}{q}
 \leq\|f^\circ\|_2^2.
\]
Creation is injective and has range $1^\perp$. Furthermore
\[
 \|f^\circ\|_2\leq\|A_\lambda f\|_2,
 \qquad f\in\Dom(A_\lambda),
\]
with equality for $f=p_1$.
\end{proposition}
\begin{proof}
Parseval proves convergence of the displayed series and its norm
identity. The finite sums belong to the polynomial graph core; their
creation images converge to $f^\circ$. Closedness of the original
creation graph therefore proves the claimed full-domain identity.
If $u=\sum u_np_n$ belongs to that graph, its creation coefficients are
$\sqrt{n+1}u_n$ in mode $n+1$. Thus $A_\lambda^*u=0$ forces every
$u_n=0$, and every creation image has zero constant coefficient.
The recovery identity proves the converse range inclusion and uniqueness
of the preimage. The maximal annihilation graph gives
\[
 \|A_\lambda f\|_2^2
 =\sum_{q\geq1}q|\langle p_q,f\rangle|^2
 \geq\sum_{q\geq1}|\langle p_q,f\rangle|^2
 =\|f^\circ\|_2^2.
\]
Only the first mode is needed for equality. Finally truncate recovery
to input modes $1,\ldots,N$. This has finite rank, and the omitted
operator has norm at most $(N+1)^{-1/2}$ by the same coefficient
identity. Thus $C_\lambda$ is an operator-norm limit of finite-rank
operators and is compact. In fact the bound is attained on $p_{N+1}$.
\end{proof}

Now fix a physical intensity change $\mu=\lambda e^\theta$, let
$V=V_{\lambda,\theta}:H_\mu\to H_\lambda$ be the root multiplication
unitary, and put $a=\sqrt\mu-\sqrt\lambda$ and $r=V1$.
The vector $r$ has norm one, but its scalar mean need not be one.
Define
\[
 P_rf=f-\langle r,f\rangle r,\qquad C_r=VC_\mu V^{-1}.
\]
\begin{corollary}\label{cor:coherent-recovery}
For every $f\in H_\lambda$,
\[
 (A_\lambda^*-a)C_rf=P_rf,\qquad
 \|C_rf\|_2\leq\|P_rf\|_2.
\]
The shifted creation graph has range $r^\perp$ and zero kernel;
$C_rf$ is its unique preimage of $P_rf$. The shifted annihilation
graph has kernel $\R r$ and satisfies
\[
 \|P_rf\|_2\leq\|(A_\lambda-a)f\|_2,
 \qquad f\in\Dom(A_\lambda).
\]
The operator $C_r$ is compact, with the same finite-rank tail bounds.
\end{corollary}
\begin{proof}
The full-domain gauge equalities proved earlier give
$(A_\lambda-a)V=VA_\mu$ and
$(A_\lambda^*-a)V=VA_\mu^*$.
The orthogonal projection onto constants in $H_\mu$ becomes projection
onto $\R r$ after conjugation. Conjugate every assertion of
Proposition~\ref{prop:recovery} by $V$. For the annihilation kernel,
its coefficient energy is zero exactly when all positive modes vanish.
Conjugating the finite-rank approximants proves compactness and the
tail bounds. Addition of the bounded scalar operator $-aI$ leaves the
respective maximal operator domains unchanged.
\end{proof}

Under $U_\lambda$, the vacuum $r$ becomes $g_a(x)=e^{ax-a^2}$.
Consequently the Gaussian operator $\delta-a$ has exact range
$g_a^\perp$, its inverse on that range is compact, and
$\ker(D-a)=\R g_a$. These are one-coordinate statements, obtained
by the same unitary conjugation. They do not assert injectivity of
divergence on a general Gaussian vector-field space. At fixed
$\lambda$, physical tilts give $a>-\sqrt\lambda$; varying $\lambda$
realizes every real Gaussian displacement. The vacuum projection,
not subtraction of the ordinary scalar mean, is the displaced
centering operation.

\section{Two intensity regimes beyond the Hilbert norm}
The preceding sharp Hilbert estimates do not yield a lower bound for
the ordinary $L^p$ norm at other exponents. This can be checked on a
single actual compensated count, without a process construction.
Let $N_r$ have Poisson law of intensity $r>0$, and put
\[
 Z_r=N_r-r,\quad b(r)=\|Z_r\|_{4/3},\quad
 s(r)=\min\{r^{1/2},r^{3/4}\}.
\]
\begin{proposition}\label{prop:two-regime}
For every $r>0$,
\[
 \tfrac14\min\{r^2,r^3\}\leq b(r)^4
 \leq4\min\{r^2,r^3\},\qquad
 \frac{s(r)}{\sqrt2}\leq b(r)\leq\sqrt2\,s(r).
\]
No positive constant $c$ satisfies
\[
 c\|1\|_{L^{4/3}(r p_r)}\leq\|Z_r\|_{L^{4/3}(p_r)}
 \quad\hbox{for every }r>0.
\]
Here $Z_r$ is precisely the original, unnormalized Poisson divergence
of the constant field $1$.
\end{proposition}
\begin{proof}
The generating function $\E t^{N_r}=e^{r(t-1)}$ gives factorial
moments $\E(N_r)_j=r^j$. Expanding ordinary powers in falling
factorials gives
\[
 \E Z_r^2=r,\qquad \E Z_r^4=3r^2+r.
\]
All moments are finite. H\"older with exponents $4/3$ and $4$, applied
to $|Z_r|\,|Z_r|$, yields
\[
 r\leq b(r)(3r^2+r)^{1/4},\qquad
 b(r)^4\geq\frac{r^3}{3r+1}.
\]
For $r\leq1$ the denominator is at most $4$; for $r\geq1$ it is
at most $4r$. This proves the lower estimate. Since the law is a
probability, monotonicity of norms gives $b(r)\leq\sqrt r$.
For the other upper bound, H\"older applied to
$|Z_r|^{2/3}|Z_r|^{2/3}$ with exponents $3/2$ and $3$ gives
\[
 \|Z_r\|_{4/3}\leq\|Z_r\|_1^{1/2}\|Z_r\|_2^{1/2}.
\]
The elementary inequality $\E|N_r-r|\leq2r$ then gives
$b(r)\leq\sqrt2\,r^{3/4}$. Combining both upper bounds proves
the assertion; taking fourth roots gives the stated comparison with
$s(r)$.

The original field space has control measure $r p_r$ and its
unnormalized divergence is
$\delta_r u(k)=k u(k-1)-r u(k)$.
For $u=1$ this equals $Z_r$. This field belongs to the original
maximal $L^2$ divergence graph: indeed the scalar unitary field
normalization maps it to $\sqrt r\,p_0$, and
$A_r^*(\sqrt r\,p_0)=\sqrt r\,p_1=Z_r$.
On the other hand
$\|1\|_{L^{4/3}(r p_r)}=r^{3/4}$.
Choose $r=n^4$ with integers $n\geq1$. The sharper probability-norm
upper bound and the lower estimate give
\[
 \frac{1}{\sqrt2\,n}\leq
 \frac{\|\delta_{n^4}1\|_{4/3}}
 {\|1\|_{L^{4/3}(n^4p_{n^4})}}\leq\frac1n.
\]
The ratio tends to zero, excluding every proposed $c>0$.
\end{proof}

The two scales have different meanings: $r^{3/4}$ describes rare
counts for $r\leq1$, while $r^{1/2}$ describes the large-intensity
regime. Only the latter regime is needed for the displayed failure of
a plain $L^{4/3}$ lower bound. There is no contradiction with the
$L^2$ estimate, where the same constant input has norm $\sqrt r$.
Every estimate transfers to a count on another probability space
having this Poisson law, since all expressions depend only on that
one law; no independence of several counts is required.

This elementary diagnostic belongs to established Poisson integration
theory. Dirksen's Theorem 1.1 \cite{Dirksen} gives much stronger
maximal stochastic integral estimates with sums and intersections of
square-function and $L^p$-type spaces. For a scalar deterministic
indicator of control mass $r$ and exponent $4/3$, its component
scales are $\sqrt r$ and $r^{3/4}$; the sum-space norm selects their
minimum up to constants. The proposition proves the count comparison
directly. It does not construct a general random measure or claim
that a single plain $L^{4/3}$ norm is the appropriate integrand space.

\section{Exact provenance and mathematical scope}
All formal source references below use the frozen verified commit
\begin{center}\code{bab688bb6140cd6316c80680095a914ea0d346db}\end{center}
Its original audit receipt and unchanged source hashes accompany the
paper in the baseline-verification JSON file. The publication editor
has not rerun a research build or altered the verified library.
The accompanying \code{theorem-map.json} records exact declarations,
line numbers and public baseline links. The source groups cover three connected layers:
\begin{itemize}
\item Complete Charlier coordinates, weighted scalar geometry and exact
Poisson spectral graphs; the Gaussian Hermite adjoint graphs.
\item Half-density source gauges, Gaussian displacement from Poisson,
coherent identification, group law, continuity and physical action.
\item Tilt Hellinger geometry, operator-norm discontinuity, quartic
unboundedness, coherent entropy distortion, finite-core failure and
original Poisson jump curvature.
\end{itemize}
The recovery and shifted Poincar\'e results are scalar normalizations of the original weighted-field graphs. The two-regime count estimates use actual Poisson laws. Compactness and the Gaussian range statements above are also proved directly by coefficients and unitary conjugation.

The finite-core entropy proof specializes the verified cutoff argument to
an explicit coherent root. Generating functions, weak-distributional
domains and the dense coherent-span identification reorganize the exact
formal graph results into human proofs; every displayed line is not a
separate formal declaration. An accompanying standard-library Python
verifier checks the Gamma-two decomposition by exact rational polynomial
arithmetic, including both boundary cases; it is a transcription check.
The positivity witness is an elementary
analytic deduction, distinguished in the theorem map.

The setting is one Poisson count and one Gaussian coordinate. The results
do not transfer ordinary nonlinear equations, a Markov order or general
entropy comparisons. This is a classical formalization companion;
novelty certification and external submission remain separate questions.

\begin{thebibliography}{9}
\bibitem{VT} A. M. Vershik and N. V. Tsilevich,
\emph{Fock factorizations, and decompositions of the $L^2$ spaces over
general L\'evy processes}, Russian Mathematical Surveys \textbf{58}
(2003), 427--472.
\href{https://arxiv.org/abs/math/0304282}{arXiv:math/0304282};
\href{https://doi.org/10.1070/RM2003v058n03ABEH000627}{doi:10.1070/RM2003v058n03ABEH000627}.
\bibitem{ALRS} S. Aryan, P. L\'opez-Rivera and Y. Shenfeld,
\emph{Improvement of Wu's logarithmic Sobolev inequality via the
Poisson--F\"ollmer process}, arXiv:2410.06117v2 (2025).
\href{https://arxiv.org/abs/2410.06117v2}{Primary manuscript},
especially the introduction, Theorem 1.1 and Section 4.
\bibitem{WZ} F. Weber and R. Zacher,
\emph{The entropy method under curvature-dimension conditions in the
spirit of Bakry--\'Emery in the discrete setting of Markov chains},
Journal of Functional Analysis \textbf{281} (2021), 109061.
\href{https://arxiv.org/abs/2007.01264}{arXiv:2007.01264}, Example 5.13.
\bibitem{Dirksen} S. Dirksen,
\emph{It\^o isomorphisms for $L^p$-valued Poisson stochastic integrals},
Annals of Probability \textbf{42} (2014), 2595--2643.
\href{https://doi.org/10.1214/13-AOP906}{doi:10.1214/13-AOP906}.
Theorem 1.1 and the norms immediately preceding it;
\href{https://research-portal.uu.nl/ws/files/107833782/13_AOP906.pdf}{primary publisher manuscript}.
\end{thebibliography}
\end{document}
